\noindent\textbf{研究要点} % qa-allow-bold (front-matter heading)
\begin{itemize}[leftmargin=1.6em]
  \item 跨期平均再生数的误差含环境方差与报告异质性两个不随规模下降的分量。
  \item 环境方差的估计精度由有效期数决定，扩大单期规模不能替代积累观测期。
  \item 环境正相关时，预算阈值决定监测应设少而大还是多而小的观测期。
  \item 可由填报医院比例观测的报告成分不到流感住院截距超出量的 1\%。
  \item 7 个病原与窗口组合中 4 个部分复现，3 个只呈方向一致，没有组合复现。
\end{itemize}
\vspace{0.8em}
